\documentclass[11pt]{article}
\usepackage[margin=1in]{geometry}
\usepackage{amsmath,amssymb,amsthm}
\usepackage{graphicx}
\usepackage{hyperref}

\newcommand{\M}{\mathcal M}
\newcommand{\h}{\mathsf h}
\newcommand{\bmo}{\mathsf{bmo}}
\newcommand{\E}{\mathcal E}

\newtheorem{theorem}{Theorem}
\newtheorem{lemma}{Lemma}

\title{A \(p\)-Lorentz-core reverse inclusion for the endpoint complex\\
noncommutative martingale Hardy--BMO problem}
\author{Automatic Math Discovery Run \texttt{fa\_banach\_001}}
\date{June 16, 2026}

\begin{document}
\maketitle

\section*{Status}

This is an unconditional partial result for the endpoint complex interpolation
problem in Randrianantoanina's paper \cite{Ran2021}. It proves that the
\(p\)-Lorentz Hardy core \(\h_{r,p}^c(\M)\) lies in the endpoint complex interpolation
space. It does not prove that all of \(\h_r^c(\M)\) lies there.

\section*{Source problem}

The source paper proves the finite-index complex interpolation theorem
\[
[\h_p^c(\M),\h_q^c(\M)]_\theta=\h_r^c(\M),
\qquad
  \frac1r=\frac{1-\theta}{p}+\frac{\theta}{q},
  \qquad 0<p,q<\infty.
\]
It then asks whether the corresponding endpoint identity remains valid when
\(0<p<1\) and the endpoint couple
\((\h_p^c(\M),\bmo^c(\M))\) is used.

\begin{center}
\includegraphics[width=\textwidth]{figures/open_problem_crop.png}
\end{center}

\noindent
The full desired endpoint identity is
\[
[\h_p^c(\M),\bmo^c(\M)]_\theta=\h_r^c(\M),
\qquad
\frac1r=\frac{1-\theta}{p}.
\]

\section*{Idea of the proof}

The Cuculescu decomposition in \cite{Ran2022} proves the real endpoint scale
\[
(\h_p^c,\h_\infty^c)_{\theta,\gamma}=\h_{r,\gamma}^c,
\qquad \frac1r=\frac{1-\theta}{p}.
\]
Choosing \(\gamma=p\) gives \(\h_{r,p}^c=(\h_p^c,\h_\infty^c)_{\theta,p}\).
The couple \((\h_p^c,\h_\infty^c)\) is naturally \(p\)-Banach: the
\(\h_p^c\) quasi-norm is \(p\)-subadditive and the \(\h_\infty^c\) norm is
also \(p\)-subadditive because \(0<p<1\). For a \(p\)-Banach couple, the real
space with summability index \(p\) embeds into the complex interpolation space
by the usual analytic series construction. Thus
\[
\h_{r,p}^c\hookrightarrow[\h_p^c,\h_\infty^c]_\theta.
\]
The endpoint inclusion \(\h_\infty^c\hookrightarrow\bmo^c\) and complex
functoriality pass this into the \(bmo^c\)-endpoint couple. The already
packaged one-sided Wolff argument gives the opposite-side bound into
\(\h_r^c\).

\begin{lemma}[real \((\theta,\rho)\) embeds into complex interpolation]
Let \((A_0,A_1)\) be a compatible couple whose endpoint quasi-norms are
\(\rho\)-subadditive for some \(0<\rho\le 1\), and let \(0<\theta<1\). Then
\[
(A_0,A_1)_{\theta,\rho}\hookrightarrow [A_0,A_1]_\theta
\]
continuously.
\end{lemma}

\begin{proof}
By the discrete \(J\)-method characterization of
\((A_0,A_1)_{\theta,\rho}\), every
\(x\) in this real interpolation space admits, up to an arbitrarily small
loss in the norm, a representation
\[
x=\sum_{n\in\mathbb Z} x_n,\qquad x_n\in A_0\cap A_1,
\]
with convergence in \(A_0+A_1\) and
\[
\sum_{n\in\mathbb Z} \big(2^{-n\theta}
J(2^n,x_n;A_0,A_1)\big)^\rho<\infty,
\qquad
J(t,u)=\max\{\|u\|_{A_0},t\|u\|_{A_1}\}.
\]
For finite partial sums define
\[
F_N(z)=\sum_{|n|\le N}2^{n(z-\theta)}x_n.
\]
Then \(F_N(\theta)=\sum_{|n|\le N}x_n\). On the boundary \(\Re z=0\),
\[
\|2^{n(z-\theta)}x_n\|_{A_0}
=2^{-n\theta}\|x_n\|_{A_0}
\le 2^{-n\theta}J(2^n,x_n),
\]
and on the boundary \(\Re z=1\),
\[
\|2^{n(z-\theta)}x_n\|_{A_1}
=2^{n(1-\theta)}\|x_n\|_{A_1}
\le 2^{-n\theta}J(2^n,x_n).
\]
The \(\rho\)-summability above and the endpoint \(\rho\)-triangle inequalities
make the series converge uniformly on both boundary lines and uniformly on
compact subsets of the strip in \(A_0+A_1\). Passing to the limit produces an
admissible analytic function \(F\) with \(F(\theta)=x\), and the complex norm
is bounded by a constant times the real \((\theta,\rho)\)-norm.
\end{proof}

\begin{lemma}
There is a continuous embedding
\[
\h_\infty^c(\M)\hookrightarrow \bmo^c(\M).
\]
\end{lemma}

\begin{proof}
For finite martingales this follows because bounded conditioned square
function controls every conditioned BMO tail:
\[
\E_n(|x_m-x_n|^2)
  =\E_n\left(\sum_{k=n+1}^m |dx_k|^2\right)
  =\E_n\left(\sum_{k=n+1}^m \E_{k-1}(|dx_k|^2)\right)
  \le \|s_c(x)^2\|_\infty{\bf 1}.
\]
The first-coordinate term is controlled by the usual \(\h_\infty^c\) norm.
Density and completion give the endpoint embedding.
\end{proof}

\begin{theorem}[\(p\)-Lorentz-core endpoint inclusion]
Let \(0<p<1\), \(0<\theta<1\), and put
\[
\frac1r=\frac{1-\theta}{p}.
\]
Assume \(r>1\). Then
\[
\h_{r,p}^c(\M)
\hookrightarrow
[\h_p^c(\M),\h_\infty^c(\M)]_\theta
\hookrightarrow
[\h_p^c(\M),\bmo^c(\M)]_\theta
\hookrightarrow
\h_r^c(\M).
\]
\end{theorem}

\begin{proof}
By Randrianantoanina's endpoint real interpolation theorem
\cite[Proposition 2.8]{Ran2022},
\[
(\h_p^c(\M),\h_\infty^c(\M))_{\theta,\gamma}=\h_{r,\gamma}^c(\M)
\]
for every \(0<\gamma\le\infty\). Taking \(\gamma=p\) gives
\[
\h_{r,p}^c(\M)=(\h_p^c(\M),\h_\infty^c(\M))_{\theta,p}.
\]
The preceding real-to-complex lemma, applied with \(\rho=p\), yields
\[
\h_{r,p}^c(\M)\hookrightarrow
[\h_p^c(\M),\h_\infty^c(\M)]_\theta.
\]
Since \(\h_\infty^c(\M)\hookrightarrow\bmo^c(\M)\), functoriality of the complex
method gives
\[
[\h_p^c(\M),\h_\infty^c(\M)]_\theta
\hookrightarrow
[\h_p^c(\M),\bmo^c(\M)]_\theta.
\]

It remains only to recall why the final inclusion holds in the range \(r>1\).
Choose \(\eta\in(0,1)\) so that
\[
1=\frac{1-\eta}{p}+\frac{\eta}{r}.
\]
The finite-index complex theorem \cite[Theorem 3.16]{Ran2021} gives
\[
[\h_p^c(\M),\h_r^c(\M)]_\eta=\h_1^c(\M).
\]
With \(\phi=1-1/r\), the Banach endpoint theorem of
Bekjan--Chen--Perrin--Yin \cite[Theorem 4.4]{BCPY2009}, after reversing the
couple, gives
\[
[\h_1^c(\M),\bmo^c(\M)]_\phi=\h_r^c(\M).
\]
Wolff's inclusion lemma, in the inclusion form used in
\cite[Corollary 3.19]{Ran2021}, therefore implies
\[
[\h_p^c(\M),\bmo^c(\M)]_\theta\hookrightarrow\h_r^c(\M),
\]
with the usual parameter identity
\(\theta=\phi/(1-\eta+\eta\phi)\). This completes the chain.
\end{proof}

\section*{Verification notes}

The proof is symbolic and uses no computation. The new inclusion is
unconditional because every input is already a theorem: the 2022 real endpoint
interpolation theorem, the elementary \(p\)-Banach
\((\theta,p)\)-real-to-complex embedding,
the endpoint \(\h_\infty^c\subset\bmo^c\) estimate, and the previously used
one-sided Wolff inclusion.

\section*{Limitations}

This does not prove the full reverse inclusion
\[
\h_r^c(\M)\hookrightarrow[\h_p^c(\M),\bmo^c(\M)]_\theta.
\]
It only proves it on the Lorentz subspace \(\h_{r,p}^c(\M)\). In general
\(L_{r,p}\subsetneq L_r\), so this should be treated as a proper partial
result. Since \(p<1\), this is a smaller core than \(\h_{r,1}^c\). The missing
part is the passage from the quasi-Banach summability index \(p\) to the
natural index \(r\), exactly the place where a stronger analytic family or
endpoint complex theorem would be needed.

\section*{Novelty check}

Local indexes were searched for arXiv:2108.06341, \(h_p^c\), \(bmo^c\),
\(h_\infty^c\), \(h_{r,p}^c\), Lorentz core, and endpoint complex
interpolation. The existing run packets contained the one-sided \(r>1\)
inclusion and a conditional \(h_\infty^c\)-endpoint reduction, but no durable
unconditional \(p\)-Lorentz-core reverse inclusion packet. External checks in the
same loop found the source paper \cite{Ran2021} and the real endpoint paper
\cite{Ran2022}; neither states the above \(p\)-Lorentz-core complex consequence as
the answer to the source endpoint problem.

\section*{Human review recommendation}

Review as an unconditional partial theorem. The main checks are the
p-Banach real-to-complex embedding
\((A_0,A_1)_{\theta,p}\hookrightarrow[A_0,A_1]_\theta\), the use of
\cite{Ran2022} with Lorentz parameter \(p\), and the endpoint embedding
\(\h_\infty^c\hookrightarrow\bmo^c\). The packet should not be read as claiming
the full reverse inclusion for \(\h_r^c\).

\begin{thebibliography}{9}

\bibitem{Ran2021}
N.~Randrianantoanina,
\emph{Interpolation between noncommutative martingale Hardy and BMO spaces:
the case \(0<p<1\)},
arXiv:2108.06341.

\bibitem{Ran2022}
N.~Randrianantoanina,
\emph{P. Jones' interpolation theorem for noncommutative martingale Hardy
spaces},
arXiv:2212.08714.

\bibitem{BCPY2009}
T.~N. Bekjan, Z.~Chen, M.~Perrin, and Z.~Yin,
\emph{Atomic decomposition and interpolation for Hardy spaces of
noncommutative martingales},
arXiv:0906.4437.

\bibitem{Wolff}
T.~Wolff,
\emph{A note on interpolation spaces},
Harmonic Analysis (Minneapolis, Minn., 1981), Lecture Notes in Math. 908,
Springer, 1982, 199--204.

\end{thebibliography}

\end{document}
